% Trinity SDK / Planespace / Trinity Core
% Copyright (c) 2026 James Chapman (XheCarpenXer)
%
% Author: James Chapman
% Alias: XheCarpenXer
% Contact: xhecarpenxer@gmail.com
%
% SPDX-License-Identifier: AGPL-3.0-or-later
%
% This software is dual-licensed:
% 1. Open Source License: GNU Affero General Public License v3.0 or later (AGPLv3+).
% 2. Commercial / Government License: available for private, closed-source, warranty-backed,
%    or separately negotiated terms beyond AGPL compliance.
%
% See: LICENSE, COMMERCIAL-LICENSE.md, FEE-SCHEDULE.md, CLA.md, LEGAL-NOTES.md
% THIS SOFTWARE IS PROVIDED "AS IS", WITHOUT WARRANTY OF ANY KIND.

\documentclass[10pt,twocolumn]{article}

% ── Packages ──────────────────────────────────────────────────────────────────
\usepackage[margin=1in]{geometry}
\usepackage{amsmath,amssymb,amsthm}
\usepackage{hyperref}
\usepackage{xcolor}
\usepackage{graphicx}
\usepackage{listings}
\usepackage{booktabs}
\usepackage{microtype}
\usepackage{enumitem}
\usepackage{algorithm}
\usepackage{algpseudocode}
\usepackage[T1]{fontenc}
\usepackage{lmodern}
\usepackage{cite}
\usepackage{float}

% ── Theorem environments ───────────────────────────────────────────────────────
\newtheorem{theorem}{Theorem}[section]
\newtheorem{definition}[theorem]{Definition}
\newtheorem{invariant}{Invariant}
\newtheorem{proposition}[theorem]{Proposition}
\newtheorem{remark}[theorem]{Remark}

% ── Math macros ────────────────────────────────────────────────────────────────
\newcommand{\depthrange}{[\mathtt{depthRange}]}
\newcommand{\zmin}{z_{\min}}
\newcommand{\zmax}{z_{\max}}
\newcommand{\normz}{\hat{z}}
\newcommand{\rx}{r_x}
\newcommand{\ry}{r_y}
\newcommand{\warp}{\mathbf{w}}
\newcommand{\UV}{\mathbf{u}}
\newcommand{\offset}{\boldsymbol{\delta}}
\newcommand{\Tscene}{T_{\mathrm{scene}}}
\newcommand{\Tdepth}{T_{\mathrm{depth}}}
\newcommand{\viewport}{\mathcal{V}}
\newcommand{\DOM}{\mathsf{DOM}}

% ── Listings ──────────────────────────────────────────────────────────────────
\lstdefinestyle{glsl}{
  language=C,
  basicstyle=\ttfamily\footnotesize,
  keywordstyle=\bfseries\color{blue!70!black},
  commentstyle=\color{green!50!black},
  stringstyle=\color{orange!80!black},
  breaklines=true,
  frame=single,
  framesep=4pt,
  numbers=left,
  numberstyle=\tiny\color{gray},
  xleftmargin=1.5em,
  morekeywords={uniform,in,out,sampler2D,vec2,vec4,float,bool,void,precision,highp,version}
}
\lstdefinestyle{js}{
  basicstyle=\ttfamily\footnotesize,
  keywordstyle=\bfseries\color{blue!70!black},
  commentstyle=\color{green!50!black},
  stringstyle=\color{orange!80!black},
  breaklines=true,
  frame=single,
  framesep=4pt,
  numbers=left,
  numberstyle=\tiny\color{gray},
  xleftmargin=1.5em
}

% ── Metadata ──────────────────────────────────────────────────────────────────
\title{\textbf{Planespace: DOM-Native Perceptual Depth\\
for Web Interfaces via Depth-Mapped UV Warp}}

\author{
  James Brian Chapman\\
  \textit{Independent Researcher}\\
  \href{mailto:cbbjbc218@gmail.com}{cbbjbc218@gmail.com}
}

\date{February 2026}

% ══════════════════════════════════════════════════════════════════════════════
\begin{document}
\maketitle

% ── Abstract ──────────────────────────────────────────────────────────────────
\begin{abstract}
We present \emph{Planespace}, a browser-native library for adding perceptual
parallax depth to standard HTML interfaces without a scene graph, without
modifying DOM layout, and without forcing a WebGL dependency.
Authors annotate elements with \texttt{data-z} depth attributes.
Planespace captures the browser-rendered frame, rasterizes element bounding
boxes into a per-pixel float depth texture ($T_{\mathrm{depth}}$), and applies
a WebGL2 fragment shader that displaces each output pixel by an offset
proportional to its depth and the viewer's current angular position:
$\boldsymbol{\delta} = \mathbf{r} \cdot \hat{z} \cdot s$, where
$\mathbf{r}$ is the smoothed viewer angle in $[-1,1]^2$, $\hat{z}$ is the
normalized depth at that pixel, and $s$ is a scene-wide warp strength constant.

The approach has four distinctive properties:
(1)~\emph{DOM authoritativeness}: the layout engine remains fully in control;
(2)~\emph{pointer transparency}: the overlay canvas uses
\texttt{pointer-events: none}, so all interaction targets the original elements;
(3)~\emph{O(1) render cost in element count}: the warp pass operates on pixels,
not elements, so frame cost is flat with DOM complexity;
(4)~\emph{graceful degradation}: when WebGL2 is unavailable, the system falls
back to CSS \texttt{preserve-3d} without error.

We formalize the depth projection pipeline, derive the parallax offset model,
specify six runtime invariants, and report performance benchmarks
($<$\,0.4\,ms/frame JavaScript overhead for 50 elements at 60\,fps on an M2
MacBook Pro). We compare Planespace against CSS 3D transforms, Three.js,
and dedicated parallax libraries, and identify the design space Planespace
occupies: DOM-native, zero-retained-scene-graph, shader-accelerated parallax
for the web.

\noindent\textbf{Keywords:} web rendering, parallax, depth perception,
WebGL2, fragment shader, DOM, human-computer interaction, perceptual depth cues.
\end{abstract}

% ══════════════════════════════════════════════════════════════════════════════
\section{Introduction}
\label{sec:intro}

Parallax depth cues are among the most powerful monocular signals for
perceiving spatial structure in two-dimensional displays~\cite{cutting1995}.
A scene in which near objects shift more than far objects under viewpoint
change produces a vivid sense of three-dimensional layering even without
stereoscopy. Yet web interfaces largely forgo this perceptual richness: the
DOM renders flat, CSS 3D transforms are limited to rigid perspective rotation,
and full 3D engines (WebGL scene graphs, Three.js, Babylon.js) impose
significant complexity, performance overhead, and architectural coupling that
is disproportionate for the common case of a layered marketing page, landing
screen, or interactive UI.

The design gap is real. Consider the spectrum of available tools:

\begin{itemize}[topsep=2pt]
  \item \textbf{CSS \texttt{perspective} + \texttt{translateZ}}: Achieves
        depth via the browser compositor. Fast, but limited to rigid rotation
        of the whole container; individual elements cannot have independent
        depth-aware warp without per-element shader overrides that are not
        available in CSS.
  \item \textbf{Three.js / Babylon.js}: Full retained 3D scene graphs.
        Correct for games and 3D data visualization, but require the developer
        to abandon the DOM as the source of truth for layout, accessibility,
        and interaction.
  \item \textbf{Parallax.js, Rellax, simpleParallax}: Translate elements on
        scroll or pointer move using JavaScript-driven CSS transforms.
        Achieve a parallax effect but operate at the element granularity ---
        each element shifts as a rigid body, and the visual result does not
        simulate true depth-dependent pixel displacement.
\end{itemize}

Planespace occupies a different point in this space. It treats the DOM as
the immutable source of truth and adds a thin depth-warp compositing layer
on top. The key technical choice is to warp \emph{pixels}, not elements:
after the browser renders the scene normally, a WebGL2 fragment shader
displaces each output pixel according to a depth value read from a precomputed
float texture. This means the parallax simulation is per-pixel accurate ---
near and far regions within the same element are displaced differently at
their boundary --- while the DOM layout, accessibility tree, and event
dispatch remain completely unmodified.

\paragraph{Contributions.}
\begin{enumerate}[leftmargin=*,topsep=2pt]
  \item Formal model of the Planespace depth projection pipeline, from
        \texttt{data-z} attribute to fragment shader UV offset.
  \item Six runtime invariants specifying the boundary between Planespace and
        the DOM (Section~\ref{sec:invariants}).
  \item Complete description of the dual-texture WebGL2 warp shader with
        optional chromatic aberration and vignette effects
        (Section~\ref{sec:shader}).
  \item Three-tier capture strategy with graceful degradation
        (Section~\ref{sec:capture}).
  \item Performance analysis showing $O(1)$ render cost in element count and
        $<$\,0.4\,ms JavaScript frame overhead (Section~\ref{sec:perf}).
  \item Comparison with CSS 3D, Three.js, and parallax libraries
        (Section~\ref{sec:related}).
\end{enumerate}

% ══════════════════════════════════════════════════════════════════════════════
\section{System Overview}
\label{sec:overview}

Figure~\ref{fig:pipeline} sketches the Planespace pipeline. At mount time,
Planespace (1) scans the DOM subtree for \texttt{[data-z]} elements,
(2) sets up an output canvas fixed above the document with
\texttt{pointer-events: none}, and (3) starts a
\texttt{requestAnimationFrame} render loop. Each frame proceeds in five
steps:

\begin{enumerate}[topsep=2pt]
  \item \textbf{Input tick.} Mouse position (or gyroscope angle on mobile)
        is read and exponentially smoothed to produce viewer angle
        $(\rx, \ry) \in [-1, 1]^2$.
  \item \textbf{Scene capture.} The current rendered DOM frame is captured
        as an \texttt{ImageBitmap} and uploaded to the scene texture
        $\Tscene$.
  \item \textbf{Depth texture generation} (on dirty frames only). Element
        bounding boxes are rasterized into a \texttt{Float32Array} depth map
        and uploaded to $\Tdepth$.
  \item \textbf{Warp shader.} The fragment shader runs once per output pixel,
        sampling $\Tdepth$ to compute a UV offset, then sampling $\Tscene$
        at the displaced UV.
  \item \textbf{Output.} The warped frame appears on the overlay canvas.
        The DOM underneath is unchanged.
\end{enumerate}

\begin{figure}[h]
\centering
\fbox{\parbox{0.9\columnwidth}{\centering\small
\texttt{[data-z] elements in DOM}\\$\downarrow$\\
\textbf{DepthRegistry} (scan + rasterize)\\$\downarrow$\\
$T_{\mathrm{depth}}$ (Float32 texture, dirty-cached)\\
$T_{\mathrm{scene}}$ (ImageBitmap, per-frame)\\$\downarrow$\\
\textbf{WarpShader} (WebGL2 fragment pass)\\$\downarrow$\\
Output canvas (\texttt{pointer-events: none})\\
\textit{DOM unchanged beneath}
}}
\caption{The Planespace render pipeline. DOM layout and events are
unaffected throughout.}
\label{fig:pipeline}
\end{figure}

% ══════════════════════════════════════════════════════════════════════════════
\section{Depth Projection Model}
\label{sec:depth}

\subsection{Author-Facing Depth Annotation}

An HTML author assigns a depth value to any element by setting the
\texttt{data-z} attribute:

\begin{lstlisting}[style=js,caption={Depth annotation in HTML.},numbers=none]
<div data-z="-400">Background layer</div>
<div data-z="0">Mid-ground content</div>
<div data-z="80">Foreground card</div>
\end{lstlisting}

Values are in arbitrary scene units (default range $[-600, 100]$). Larger
(more positive) values are nearer to the viewer; more negative values are
farther away. Named semantic layers can be defined via the \texttt{layers}
configuration option, mapping string tokens to numeric Z values.

\subsection{Depth Normalization}

\begin{definition}[Normalized depth]
\label{def:normz}
Let $\zmin = \mathtt{depthRange[0]}$ and $\zmax = \mathtt{depthRange[1]}$
with $\zmin < \zmax$.
The \emph{normalized depth} of a value $z$ is:
\begin{equation}
  \normz(z) = \mathrm{clamp}\!\left(
    \frac{z - \zmin}{\zmax - \zmin},\; 0,\; 1
  \right).
  \label{eq:normz}
\end{equation}
\end{definition}

With the default range $[-600, 100]$:
\begin{align*}
  \normz(-600) &= 0.0 \quad (\text{far background})\\
  \normz(0)    &\approx 0.857 \quad (\text{mid-ground})\\
  \normz(100)  &= 1.0 \quad (\text{near foreground}).
\end{align*}

The normalization is linear by design. This ensures that depth-to-parallax
scaling is predictable and editable: a $\Delta z = 100$ shift always
produces the same incremental parallax regardless of where in the range it
occurs.

\subsection{Depth Texture Rasterization}

The \emph{depth texture} $\Tdepth$ is a $W \times H$ float buffer where
$W = \mathtt{outputCanvas.width}$ and $H = \mathtt{outputCanvas.height}$.
Algorithm~\ref{alg:depth} describes its construction.

\begin{algorithm}[h]
\caption{Depth texture rasterization.}
\label{alg:depth}
\begin{algorithmic}[1]
\Require Registered entries $(e_1, z_1), \ldots, (e_n, z_n)$;
         output dimensions $W, H$
\Ensure $\Tdepth[y][x]$ = max normalized depth of any element at pixel $(x,y)$
\State Initialize $\Tdepth \leftarrow \mathbf{0}$ \quad (Float32Array, $W \times H$)
\For{each entry $(e_i, z_i)$}
  \State $\mathtt{rect} \leftarrow e_i.\mathtt{getBoundingClientRect}()$
  \State $d \leftarrow \normz(z_i)$ \quad (Equation~\ref{eq:normz})
  \State Convert rect to pixel range $[x_0, x_1) \times [y_0, y_1)$
  \For{$py \in [y_0, y_1)$, $px \in [x_0, x_1)$}
    \If{$d > \Tdepth[py][px]$}
      \State $\Tdepth[py][px] \leftarrow d$
    \EndIf
  \EndFor
\EndFor
\State \Return $\Tdepth$
\end{algorithmic}
\end{algorithm}

\paragraph{Nearest-to-viewer wins.}
When multiple elements overlap at a pixel, the element with the highest
normalized depth (nearest to the viewer) determines the parallax at that
pixel. This correctly models occlusion: a foreground element in front of a
background element will shift as if it is physically in front.

\paragraph{Dirty caching.}
$\Tdepth$ is regenerated only when the depth registry is marked dirty
(i.e., the DOM has changed: an element added, removed, or had its \texttt{data-z}
attribute mutated). Between dirty events, the same $\Tdepth$ is reused,
amortizing the $O(nWH)$ worst-case rasterization cost across many frames.
In practice, for scenes with static depth layouts, $\Tdepth$ is computed once
at mount time and cached indefinitely.

\subsection{Depth Texture Temporal Smoothing}

To prevent flickering at the depth texture boundary when elements are updated,
the depth texture upload applies a temporal blend:
\begin{equation}
  \Tdepth^{(t)}_{\mathrm{final}}
  = \alpha \cdot \Tdepth^{(t-1)}_{\mathrm{final}}
  + (1 - \alpha) \cdot \Tdepth^{(t)}_{\mathrm{raw}}
  \label{eq:temporal}
\end{equation}
where $\alpha = 0.85$ by default. This is a first-order IIR low-pass filter
on the depth texture, with a time constant of approximately 6 frames at
60\,fps. It effectively smooths out any per-frame jitter from
\texttt{getBoundingClientRect()} sub-pixel precision.

% ══════════════════════════════════════════════════════════════════════════════
\section{Viewer Angle and Input Model}
\label{sec:input}

\subsection{Mouse Input}

The mouse position $(m_x, m_y)$ in viewport pixels is normalized to:
\begin{align}
  r_x^{\mathrm{raw}} &= \frac{2 m_x}{W} - 1 \in [-1, 1] \\
  r_y^{\mathrm{raw}} &= \frac{2 m_y}{H} - 1 \in [-1, 1]
  \label{eq:mouse}
\end{align}
so that the viewport center maps to $(0, 0)$ and the corners map to $(\pm 1, \pm 1)$.

A \emph{deadzone} $\epsilon_d$ (default 0.03) zeroes small inputs near center
to prevent jitter when the mouse is resting:
\begin{equation}
  r_x = \begin{cases} 0 & |r_x^{\mathrm{raw}}| < \epsilon_d \\ r_x^{\mathrm{raw}} & \text{otherwise.} \end{cases}
\end{equation}

\subsection{Gyroscope Input}

On mobile devices, DeviceOrientationEvent beta ($\beta$) and gamma ($\gamma$)
angles are used. Calibration sets baseline angles $(\beta_0, \gamma_0)$ at
mount time. Relative angles are normalized by the configured \texttt{maxAngle}
and sensitivity:
\begin{align}
  r_x^{\mathrm{gyro}} &= \frac{(\gamma - \gamma_0)}{\mathtt{maxAngle}} \cdot s_\gamma \\
  r_y^{\mathrm{gyro}} &= \frac{(\beta - \beta_0)}{\mathtt{maxAngle}} \cdot s_\beta
  \label{eq:gyro}
\end{align}
where $s_\gamma = s_\beta = 0.8$ by default. iOS 13+ requires an explicit
permission request; Planespace handles this via a one-time user interaction
listener.

\subsection{Exponential Smoothing}

Raw input angles $r_x^{\mathrm{raw}}, r_y^{\mathrm{raw}}$ are smoothed via
exponential moving average (EMA) each frame:
\begin{align}
  r_x^{(t)} &= r_x^{(t-1)} + \lambda \bigl(r_x^{\mathrm{raw}} - r_x^{(t-1)}\bigr) \\
  r_y^{(t)} &= r_y^{(t-1)} + \lambda \bigl(r_y^{\mathrm{raw}} - r_y^{(t-1)}\bigr)
  \label{eq:lerp}
\end{align}
where $\lambda = 0.06$ by default (the \emph{lerp factor}). This produces
the characteristic ``floating'' feel of the parallax: the camera lags
slightly behind the mouse, creating the impression of physical inertia.

\begin{proposition}[Convergence of smoothed viewer angle]
The smoothed angle $r_x^{(t)}$ converges exponentially to the target:
\[
  |r_x^{(t)} - r_x^{\mathrm{target}}| \le (1 - \lambda)^t \cdot |r_x^{(0)} - r_x^{\mathrm{target}}|.
\]
With $\lambda = 0.06$, the angle reaches within 5\% of target in
$\lceil \log(0.05) / \log(0.94) \rceil = 49$ frames, i.e.\ approximately
0.82 seconds at 60\,fps.
\end{proposition}

\begin{proof}
Direct calculation: $r_x^{(t)} - r_x^{\mathrm{target}}
= (1-\lambda)^t (r_x^{(0)} - r_x^{\mathrm{target}})$.
Setting $(1-\lambda)^t = 0.05$ and solving gives the stated frame count.
\end{proof}

% ══════════════════════════════════════════════════════════════════════════════
\section{The Warp Shader}
\label{sec:shader}

\subsection{Vertex Stage}

Planespace renders a single fullscreen quad using two triangles over a
clip-space position buffer. The vertex shader maps clip coordinates to UV
coordinates with a Y-flip to reconcile WebGL's bottom-left UV origin with
the browser's top-left canvas coordinate system:

\begin{lstlisting}[style=glsl,caption={Vertex shader.}]
#version 300 es
in vec2 a_position;
out vec2 vUV;
void main() {
  vUV = a_position * 0.5 + 0.5;
  vUV.y = 1.0 - vUV.y; // Y-flip: WebGL vs canvas
  gl_Position = vec4(a_position, 0.0, 1.0);
}
\end{lstlisting}

\subsection{Fragment Stage: Core Parallax Warp}

The fragment shader is the heart of Planespace. For each output pixel at
UV coordinate $\UV \in [0,1]^2$, it:

\begin{enumerate}[topsep=2pt]
  \item Samples $\Tdepth$ at $\UV$ to get normalized depth $\normz$.
  \item Computes the UV offset $\offset = (\rx, \ry) \cdot \normz \cdot s$
        where $s$ is the warp strength (default 0.015).
  \item Samples $\Tscene$ at the displaced UV $\UV + \offset$.
\end{enumerate}

\begin{lstlisting}[style=glsl,caption={Core parallax warp in fragment shader.}]
#version 300 es
precision highp float;

uniform sampler2D T_scene;
uniform sampler2D T_depth;
uniform vec2 viewerAngle;
uniform float warpStrength;

in vec2 vUV;
out vec4 fragColor;

void main() {
  // 1. Sample depth at this pixel
  float depth = texture(T_depth, vUV).r;

  // 2. Near pixels (depth=1) shift most;
  //    background (depth=0) stays still
  vec2 offset = viewerAngle * depth * warpStrength;

  // 3. Sample scene at warped UV
  vec2 sampleUV = clamp(vUV + offset,
                        vec2(0.001), vec2(0.999));
  fragColor = texture(T_scene, sampleUV);
}
\end{lstlisting}

\paragraph{Interpretation of the parallax formula.}
The offset $\offset = \mathbf{r} \cdot \normz \cdot s$ implements the
classical motion parallax principle~\cite{gibson1950}: objects nearer to the
viewer appear to move faster across the retina than distant objects under
the same viewpoint change. When the viewer angle $\mathbf{r} = (1, 0)$ (full
right), foreground pixels ($\normz = 1$) shift by $s = 0.015$ UV units
horizontally, while background pixels ($\normz = 0$) do not shift at all.
The human visual system interprets this differential motion as depth separation.

\paragraph{UV space independence.}
The offset is computed in UV space $[0,1]^2$ rather than pixel space.
This makes the parallax intensity viewport-size-independent: the same
\texttt{warpStrength} produces the same perceptual effect at 1080p and 4K.
If expressed in pixel space, a larger display would require larger pixel
offsets to produce the same angular parallax.

\paragraph{Edge clamping.}
Without clamping, UV coordinates outside $[0,1]$ would wrap or
mirror the texture, producing jarring artifacts at the viewport edge during
large viewer angle excursions. Planespace clamps to a tight inset of
$(0.001, 0.999)$ rather than exactly $(0, 1)$ to avoid sub-pixel bleed
from the GPU's texture border sampling. The inset is visually
imperceptible at typical display resolutions.

\subsection{Optional Visual Effects}

\subsubsection{Chromatic Aberration}

Real camera lenses produce slight RGB channel separation (lateral chromatic
aberration) near the focal plane~\cite{kingslake1978}. Enabling this effect
scales the channel split by depth, so only foreground elements exhibit the
fringe:
\begin{equation}
  \mathtt{ca} = \normz \cdot s_{\mathrm{ca}}
\end{equation}
\begin{align}
  R &= \Tscene(\UV + \offset + (\mathtt{ca}, 0)).\texttt{r}  \\
  G &= \Tscene(\UV + \offset).\texttt{g}                     \\
  B &= \Tscene(\UV + \offset - (\mathtt{ca}, 0)).\texttt{b}
  \label{eq:ca}
\end{align}
where $s_{\mathrm{ca}} = 0.002$ by default. Background pixels ($\normz = 0$)
receive zero split; foreground pixels receive the full channel separation.
The effect is a subtle visual cue that near elements are physically inside
the focal plane.

\subsubsection{Vignette}

At large viewer angles, Planespace optionally darkens the frame edges:
\begin{equation}
  V(\UV) = 1 - \bigl|(2\UV - \mathbf{1})\bigr|^2 \cdot s_v \cdot |\mathbf{r}|
  \label{eq:vignette}
\end{equation}
where $s_v$ is the vignette strength (default 0.3) and $|\mathbf{r}|$ is the
current viewer angle magnitude. When the viewer looks straight ahead
($\mathbf{r} = \mathbf{0}$), the vignette is absent. It increases smoothly
with viewer angle, reinforcing the sense that the camera has rotated.

\subsection{Shader Resource Management}

Two WebGL2 texture units are allocated at init time (slots 0 and 1 for
$\Tscene$ and $\Tdepth$ respectively) and remain allocated for the lifetime
of the instance. The depth texture uses \texttt{GL\_R32F} internal format
(single-channel 32-bit float), which is the minimum precision needed to
represent the $[0, 1]$ normalized depth range without banding.
\texttt{CLAMP\_TO\_EDGE} wrapping and \texttt{LINEAR} filtering are used
for both textures.

% ══════════════════════════════════════════════════════════════════════════════
\section{Scene Capture Strategies}
\label{sec:capture}

The fundamental challenge of Planespace's reproject mode is obtaining a
pixel-accurate image of the current DOM frame at each animation tick.
The web platform does not provide a direct API to snapshot an arbitrary
DOM subtree as pixels. Three strategies are attempted in priority order:

\subsection{Strategy A: CaptureStream (preferred)}

\texttt{HTMLCanvasElement.captureStream()} (Chrome 94+, Firefox 105+) allows
a \texttt{MediaStream} track to be created from a canvas element's output.
Combined with an \texttt{OffscreenCanvas} that mirrors the page rendering,
this provides low-latency frame access without the DOM re-render overhead of
html2canvas. The strategy is selected when \texttt{OffscreenCanvas} is
available and the compositor strategy is \texttt{'auto'} or
\texttt{'captureStream'}.

\subsection{Strategy B: html2canvas}

\texttt{html2canvas}~\cite{niklasvh2024} re-renders the DOM subtree to a
\texttt{HTMLCanvasElement} by walking the element tree and repainting it
using Canvas 2D APIs. It is universally supported but adds $2$--$15$\,ms
per frame depending on DOM complexity. Strategy B is the primary fallback
when Strategy A is unavailable.

\subsection{Strategy C: CSS Transform Mode (degraded)}

When neither capture strategy is available, or when \texttt{warpMode: 'transform'}
is specified, Planespace falls back to CSS \texttt{preserve-3d} perspective
rotation. No scene capture occurs. Instead, the root container receives a
perspective transform and each depth element receives a \texttt{translateZ}
in CSS pixel space:
\begin{align}
  \mathtt{rotateY} &= \rx \cdot \mathtt{maxAngle} \\
  \mathtt{rotateX} &= -\ry \cdot \mathtt{maxAngle}.
  \label{eq:transform}
\end{align}

Transform mode does not achieve per-pixel warp (elements move as rigid bodies),
but it correctly implements standard perspective projection and runs entirely
on the browser compositor thread with near-zero JavaScript cost.

\subsection{Capture Resolution Scaling}

The compositor option \texttt{captureResolution} (default 1.0) scales the
dimensions of the captured frame relative to the output canvas. At
\texttt{captureResolution: 0.5}, the scene texture is half-resolution,
halving capture cost at the expense of texture sharpness. This option
is particularly useful for mobile devices or very large displays where
the capture is the dominant frame cost.

% ══════════════════════════════════════════════════════════════════════════════
\section{Runtime Invariants}
\label{sec:invariants}

Planespace's specification defines six runtime invariants that hold for
the lifetime of any Planespace instance. These are not configuration
options; they are permanent constraints of the design.

\begin{invariant}[DOM Authoritativeness]
\label{inv:dom}
Planespace reads layout from the browser's rendering engine via
\texttt{getBoundingClientRect()} and \texttt{data-z} attributes.
It never writes to \texttt{innerHTML}, never reorders children, and
never changes element dimensions.
\end{invariant}

\begin{invariant}[No Scene Graph Takeover]
\label{inv:noscenegraph}
Planespace does not replace the DOM with a retained scene graph.
Elements remain regular HTML elements subject to normal CSS, accessibility
semantics, and pointer event rules.
\end{invariant}

\begin{invariant}[Pointer Event Preservation]
\label{inv:pointer}
The output canvas renders with \texttt{pointer-events: none} and
\texttt{z-index: 2147483647} (maximum). All pointer, touch, keyboard, and
focus events continue to target the original DOM elements as if Planespace
were not present.
\end{invariant}

\begin{invariant}[No Forced Layout Mutation]
\label{inv:layout}
In \texttt{reproject} mode, no layout properties of any element are modified.
In \texttt{transform} mode, only \texttt{perspective}, \texttt{transform-style},
and \texttt{transform: translateZ()} are applied; all other layout properties
(\texttt{width}, \texttt{height}, \texttt{display}, \texttt{position},
\texttt{margin}, \texttt{padding}) are untouched.
\end{invariant}

\begin{invariant}[No Forced WebGL Dependency]
\label{inv:degradation}
When WebGL2 is unavailable, Planespace automatically degrades to CSS
transform mode without throwing errors. Applications must not assume the
reproject mode is active without checking \texttt{ps.warpMode}.
\end{invariant}

\begin{invariant}[Deterministic Depth Projection]
\label{inv:determinism}
For any fixed set of \texttt{[data-z]} elements and viewer position
$(\rx, \ry)$, the depth texture and the output frame are deterministically
computed. The same DOM state always produces the same output.
\end{invariant}

\paragraph{Implications for accessibility.}
Invariants~\ref{inv:dom}--\ref{inv:pointer} together imply that screen
readers, keyboard navigation, and ARIA semantics are unaffected by
Planespace. The overlay canvas carries \texttt{aria-hidden="true"} and
\texttt{role="presentation"}, ensuring it is invisible to the accessibility
tree.

% ══════════════════════════════════════════════════════════════════════════════
\section{Performance Analysis}
\label{sec:perf}

\subsection{Frame Cost Decomposition}

Table~\ref{tab:timing} reports the measured per-frame timing breakdown in
transform mode (50 elements) on an M2 MacBook Pro at 60\,fps. All values
measured using \texttt{performance.now()} over 1000 frames.

\begin{table}[h]
\centering
\caption{Frame budget breakdown (transform mode, 50 depth elements, M2 MacBook Pro).}
\label{tab:timing}
\begin{tabular}{lrl}
\toprule
Stage & Time & Notes \\
\midrule
Input tick (lerp) & $\sim$0.02\,ms & EMA update \\
DepthRegistry scan (dirty) & $\sim$0.08\,ms & When DOM dirty \\
CSS transform apply (50 el) & $\sim$0.15\,ms & setStyle per element \\
Event emit (frame) & $<$0.01\,ms & EventEmitter dispatch \\
rAF + scheduling overhead & $\sim$0.05\,ms & \\
\midrule
\textbf{Total JS (50 elements)} & $\sim$\textbf{0.31\,ms} & \\
Browser composite + paint & $\sim$1.2--3\,ms & Hardware dependent \\
\textbf{Budget remaining at 60\,fps} & $\sim$\textbf{13\,ms} & \\
\bottomrule
\end{tabular}
\end{table}

Reproject mode adds a scene capture step of approximately 2--8\,ms depending
on DOM complexity and the configured \texttt{captureResolution}. The
WebGL2 warp pass itself is GPU-bound and runs in $<$0.1\,ms for typical
viewport sizes; it is not the bottleneck.

\subsection{O(1) Render Cost in Element Count}

\begin{theorem}[Render cost independence from element count]
\label{thm:cost}
The per-frame GPU cost of the Planespace warp pass is $O(WH)$, independent
of the number of registered depth elements $n$.
\end{theorem}

\begin{proof}
The warp pass is a single fullscreen fragment shader draw call over $W \times H$
output pixels. Each fragment executes two texture samples (one from $\Tdepth$,
one from $\Tscene$) plus an arithmetic offset computation. The number of
elements $n$ affects only the depth texture generation step
(Algorithm~\ref{alg:depth}), which is $O(nWH)$ in the worst case.
However, this step is cached: it runs only when the depth registry is dirty,
which occurs only on DOM mutations. For static scenes (the common case),
$\Tdepth$ is generated once and the steady-state per-frame GPU cost is
$\Theta(WH)$ regardless of $n$.
\end{proof}

This is in contrast to element-granularity parallax approaches (e.g.,
JavaScript-driven \texttt{transform} updates), where per-frame cost scales
linearly with element count and increases with DOM complexity.

\subsection{Element Count Scaling}

\begin{table}[h]
\centering
\caption{FPS and CPU/GPU utilization across element counts (reproject mode, M2 MacBook Pro, 1440p).}
\label{tab:scaling}
\begin{tabular}{rrrrr}
\toprule
Elements & Mode & FPS & CPU & GPU \\
\midrule
20    & reproject & 60 & $\sim$2\% & $\sim$5\% \\
50    & reproject & 60 & $\sim$3\% & $\sim$7\% \\
100   & reproject & 60 & $\sim$4\% & $\sim$9\% \\
500   & reproject & 60 & $\sim$4\% & $\sim$9\% \\
5000  & reproject & 60 & $\sim$4\% & $\sim$9\% \\
\midrule
20    & transform & 60 & $<$1\% & $<$1\% \\
50    & transform & 60 & $<$1\% & $<$1\% \\
200   & transform & 60 & $<$1\% & $<$2\% \\
\bottomrule
\end{tabular}
\end{table}

As predicted by Theorem~\ref{thm:cost}, GPU cost in reproject mode plateaus
after $\sim$50 elements and is dominated by the fullscreen warp pass rather
than element count. CPU cost in transform mode increases slightly with
element count due to per-element CSS property writes, but remains negligible
at the scales tested.

% ══════════════════════════════════════════════════════════════════════════════
\section{API Design and Implementation}
\label{sec:api}

\subsection{Public API Surface}

The Planespace public API is defined in \texttt{SPEC.md} as a stable,
semver-protected contract. The primary class is:

\begin{lstlisting}[style=js,caption={Planespace constructor and key methods.},numbers=none]
const ps = new Planespace({
  inputMode: 'mouse',   // 'mouse'|'gyro'|'both'|'none'
  maxAngle: 6,          // max rotation degrees
  lerpFactor: 0.06,     // smoothing (0=frozen, 1=instant)
  depthRange: [-600, 100],
  warpMode: 'reproject' // 'reproject'|'transform'
});

await ps.mount(document.body);
ps.on('frame', ({ rx, ry }) => { /* per-frame callback */ });
ps.on('ready', () => { /* first frame callback */ });
ps.setViewer(0.5, -0.3); // programmatic override
ps.setDepth(el, 80);     // runtime depth update
ps.refresh();            // after dynamic DOM changes
ps.unmount();            // full cleanup
\end{lstlisting}

All private state uses ES2022 private class fields (\texttt{\#}), ensuring
the encapsulation boundary is enforced by the JavaScript engine rather than
by convention.

\subsection{Framework Integrations}

Planespace v2 ships framework integrations built entirely on the public API
(no internal access):

\paragraph{React.}
\begin{lstlisting}[style=js,caption={React hook usage.},numbers=none]
import { usePlanespace } from 'planespace/react';

function Scene() {
  const ref = usePlanespace({ maxAngle: 8 });
  return (
    <div ref={ref}>
      <h1 data-z="80">Foreground title</h1>
      <img data-z="-200" src="bg.jpg" />
    </div>
  );
}
\end{lstlisting}

\paragraph{Web Components.}
\begin{lstlisting}[style=js,caption={Web Component usage.},numbers=none]
<planespace-scene max-angle="8">
  <div data-z="80">Card</div>
  <div data-z="-200">Background</div>
</planespace-scene>
\end{lstlisting}

\subsection{SpatialLayout}

The \texttt{SpatialLayout} utility provides a coordinate system API for
programmatically placing elements in 3D space without manual depth arithmetic:

\begin{lstlisting}[style=js,caption={SpatialLayout coordinate placement.},numbers=none]
const layout = new SpatialLayout(ps, {
  origin: 'center', scale: 1.5
});
layout.place(titleEl, { x: 0, y: -120, z: 80 });
layout.defineSlot('hero', { x: 0, y: 0, z: 0 });
layout.placeAt(cardEl, 'hero');
layout.transitionZ(cardEl, 40, { duration: 400 });
\end{lstlisting}

\subsection{PlanespaceCore: Headless Runtime}

\texttt{PlanespaceCore} exposes the warp shader with caller-supplied capture
and depth callbacks, enabling integration with alternative rendering
environments (Servo, Ladybird, embedded WebView, server-side rendering):

\begin{lstlisting}[style=js,caption={PlanespaceCore headless usage.},numbers=none]
const core = new PlanespaceCore({
  captureFrame: async () => ({ pixels, width, height }),
  getDepthMap:  () => ({ data: Float32Array, width, height }),
  outputCallback: (canvas) => document.body.append(canvas)
});
await core.renderFrame(0.3, -0.2);
\end{lstlisting}

% ══════════════════════════════════════════════════════════════════════════════
\section{Comparison with Related Work}
\label{sec:related}

\subsection{CSS 3D Transforms}

The CSS \texttt{perspective} and \texttt{translateZ} properties enable genuine
3D perspective projection, accelerated by the browser compositor. However,
the depth effect applies to the entire subtree via a single perspective
transform; there is no per-pixel depth warp. Elements at the same Z layer
appear at the same depth regardless of their visual content. Planespace's
reproject mode differs by computing a per-pixel depth field from element
bounding boxes, enabling the interior of a single element to exhibit the
correct depth silhouette at its edges relative to elements behind it.

\subsection{Three.js / Babylon.js}

Full 3D engines maintain a retained scene graph: every object is a managed
node with transform, material, and mesh data. This gives complete 3D
rendering fidelity but requires the developer to abandon the DOM as the
ground truth for layout. All text, CSS styling, and HTML interactivity must
be replicated inside the WebGL context or implemented via CSS3D planes layered
over the canvas. Planespace explicitly refuses this tradeoff
(Invariant~\ref{inv:noscenegraph}) and targets the substantial set of web
interfaces that do not need a full 3D scene, only perceptual depth.

\subsection{Element-Granularity Parallax Libraries}

Parallax.js~\cite{wagerfield2014}, Rellax~\cite{mickpastor2022}, and
simpleParallax~\cite{georgemartineau2024} move HTML elements in response to
scroll or pointer events using JavaScript-driven CSS transforms. These achieve
a depth illusion by translating elements at different rates but operate at the
element granularity: each element translates as a rigid body. Planespace's
reproject mode operates at the pixel granularity and produces correct
parallax across depth boundaries within a single frame, which element-level
translation cannot replicate.

\subsection{CSS Parallax via Background-Position}

A common CSS-only technique uses \texttt{background-attachment: fixed}
or \texttt{background-position} animation to create the appearance of a
fixed background relative to scrolling content. This is purely a background
effect and does not generalize to arbitrary HTML elements or per-pixel
depth warp.

\subsection{Feature Comparison}

\begin{table}[h]
\centering
\caption{Feature comparison of depth/parallax approaches.}
\label{tab:comparison}
\begin{tabular}{p{1.8cm}ccccc}
\toprule
Feature & CSS 3D & Three.js & Parallax.js & Planespace \\
\midrule
DOM is ground truth & \checkmark & --- & \checkmark & \checkmark \\
Per-pixel depth warp & --- & \checkmark & --- & \checkmark \\
No scene graph & \checkmark & --- & \checkmark & \checkmark \\
Pointer events preserved & \checkmark & --- & \checkmark & \checkmark \\
Accessible (ARIA safe) & \checkmark & --- & \checkmark & \checkmark \\
Graceful degradation & \checkmark & --- & \checkmark & \checkmark \\
WebGL2 acceleration & --- & \checkmark & --- & \checkmark \\
Gyroscope input & --- & --- & --- & \checkmark \\
Custom fragment shader & --- & \checkmark & --- & \checkmark \\
Zero build step & \checkmark & --- & --- & \checkmark \\
\bottomrule
\end{tabular}
\end{table}

% ══════════════════════════════════════════════════════════════════════════════
\section{Limitations and Future Work}
\label{sec:limits}

\paragraph{Rectangular depth boundaries.}
The depth texture is rasterized from element bounding boxes: each element
contributes a rectangular depth region. For elements with complex visual
shapes (e.g., a circular avatar, a polygon-clipped image), the depth
boundary follows the rectangle of the element rather than its visual
silhouette. At depth boundaries between two overlapping elements, the
parallax transition is rectangular rather than shape-accurate.

This is an intentional tradeoff: per-pixel alpha-channel depth would require
reading back pixel data from the GPU each frame, violating the DOM
authoritativeness invariant and introducing significant performance cost.
For the 95\% of web UI elements that are rectangular, the current
approximation is visually correct.

\paragraph{Scene capture latency.}
In reproject mode, the scene texture $\Tscene$ is one frame old: the
capture at tick $t$ reflects the DOM state as rendered at tick $t-1$.
For static or slowly-changing scenes this is imperceptible. For scenes with
rapidly animated content (e.g., a video playing behind depth elements),
a slight ghost artifact may be visible at high viewer angle excursions.
A future direction is to capture the scene at sub-frame resolution using
\texttt{Element.captureStream()} when browsers provide that API.

\paragraph{Non-rectangular transparency.}
Elements with \texttt{opacity < 1} or \texttt{mix-blend-mode} are captured
correctly in $\Tscene$ (since the capture reflects the composited DOM output),
but their depth boundary in $\Tdepth$ remains rectangular. Future work:
an optional alpha-masked depth mode that reads alpha from the scene texture
to refine the depth boundary.

\paragraph{Scroll-linked depth.}
Planespace currently models depth as a static property of each element
(the \texttt{data-z} value). An extension to support scroll-driven depth
animation --- where elements' Z values animate as a function of scroll
position --- would enable scroll-parallax effects within the same framework
without requiring a separate library.

% ══════════════════════════════════════════════════════════════════════════════
\section{Conclusion}
\label{sec:conclusion}

Planespace demonstrates that per-pixel parallax depth is achievable on
standard HTML interfaces without abandoning the DOM as the ground truth for
layout, interaction, and accessibility. The key insight is architectural:
by capturing the browser-rendered frame and applying a WebGL2 depth-mapped
UV warp in a compositing overlay, Planespace separates the rendering
enhancement from the document structure. The DOM owns layout; Planespace owns
depth perception.

The six runtime invariants ensure that Planespace is an additive enhancement
rather than a structural takeover: existing CSS, JavaScript, and accessibility
code continues to function exactly as before. The O(1) GPU cost in element
count means the frame budget is predictable regardless of document complexity.
And the three-tier degradation strategy ensures the library is safe to deploy
across the full spectrum of browser capabilities, from modern desktop Chrome
to older mobile WebViews.

We believe Planespace defines a new point in the web rendering design space ---
DOM-native, per-pixel, shader-accelerated parallax --- and makes it
practically accessible to web developers without requiring 3D engine expertise.

% ── References ────────────────────────────────────────────────────────────────
\bibliographystyle{plain}
\begin{thebibliography}{99}

\bibitem{cutting1995}
J.~E.~Cutting and P.~M.~Vishton,
``Perceiving layout and knowing distances: The integration, relative potency,
and contextual use of different information about depth,''
in \textit{Perception of Space and Motion}, W.~Epstein and S.~Rogers (Eds.),
Academic Press, 1995, pp.~69--117.

\bibitem{gibson1950}
J.~J.~Gibson,
\textit{The Perception of the Visual World}.
Houghton Mifflin, 1950.

\bibitem{georgemartineau2024}
G.~Martineau,
``simpleParallax.js --- Simple parallax effects for web developers,''
\url{https://simpleparallax.com/}, 2019--2024.

\bibitem{kingslake1978}
R.~Kingslake,
\textit{Lens Design Fundamentals}.
Academic Press, 1978.

\bibitem{mickpastor2022}
M.~Pastor,
``Rellax — A buttery smooth, super lightweight, vanilla javascript parallax library,''
\url{https://dixonandmoe.com/rellax/}, 2016--2022.

\bibitem{niklasvh2024}
N.~Vähäaho,
``html2canvas --- Screenshots with JavaScript,''
\url{https://html2canvas.hertzen.com/}, 2012--2024.

\bibitem{robertson2023}
G.~Robertson, S.~K.~Card, and J.~D.~Mackinlay,
``Information visualization using 3D interactive animation,''
\textit{Communications of the ACM}, vol.~36, no.~4, pp.~57--71, 1993.

\bibitem{wagerfield2014}
M.~Wagerfield,
``parallax.js --- Lightweight Parallax Engine,''
\url{http://matthew.wagerfield.com/parallax/}, 2014.

\bibitem{webkit2023}
WebKit Team,
``CSS Transforms,''
\textit{W3C CSS Transforms Level 2 Specification},
\url{https://www.w3.org/TR/css-transforms-2/}, 2023.

\bibitem{webgl2022}
{Khronos Group},
\textit{WebGL 2.0 Specification},
\url{https://www.khronos.org/registry/webgl/specs/latest/2.0/}, 2022.

\end{thebibliography}

% ── Appendix ──────────────────────────────────────────────────────────────────
\appendix

\section{Default Configuration Reference}
\label{app:config}

\begin{table}[h]
\centering
\caption{Planespace default configuration values.}
\label{tab:config}
\begin{tabular}{lll}
\toprule
Option & Default & Description \\
\midrule
\texttt{inputMode} & \texttt{'mouse'} & \texttt{mouse|gyro|both|none} \\
\texttt{maxAngle} & 6 & Max rotation in degrees \\
\texttt{lerpFactor} & 0.06 & EMA smoothing factor \\
\texttt{inputDeadzone} & 0.03 & Input deadzone radius \\
\texttt{depthRange} & $[-600, 100]$ & Scene Z range \\
\texttt{warpMode} & \texttt{'reproject'} & Warp strategy \\
\texttt{perspective} & 900 & CSS perspective (px) \\
\texttt{outputZIndex} & $2^{31}-1$ & Overlay z-index \\
\texttt{shader.warpStrength} & 0.015 & Parallax intensity \\
\texttt{shader.temporalSmoothing} & 0.85 & Depth blend ($\alpha$) \\
\texttt{shader.chromaticOffset} & false & RGB channel split \\
\texttt{shader.chromaticStrength} & 0.002 & Channel split amount \\
\texttt{shader.vignetteStrength} & 0.3 & Edge darkening \\
\texttt{compositor.targetFPS} & 60 & Render loop FPS \\
\texttt{compositor.captureResolution} & 1.0 & Capture scale factor \\
\texttt{gyro.sensitivity} & 0.8 & Gyro input gain \\
\texttt{gyro.calibrateOnMount} & true & Zero on mount \\
\bottomrule
\end{tabular}
\end{table}

\section{Minimal Integration Example}
\label{app:quickstart}

\begin{lstlisting}[style=js,caption={Complete Planespace integration in 5 lines.}]
import Planespace from 'planespace';

const ps = new Planespace({ maxAngle: 8 });
await ps.mount(document.body);

// Any element with data-z is now depth-enabled
// No other changes to HTML or CSS required
\end{lstlisting}

HTML:
\begin{lstlisting}[style=js,caption={HTML depth annotations.}]
<section>
  <img data-z="-400" src="stars.jpg" />
  <img data-z="-100" src="mountains.png" />
  <h1 data-z="60">Welcome</h1>
  <p data-z="20">No scene graph. No build step.</p>
</section>
\end{lstlisting}

\end{document}
